\chapter{CONCLUSIONS AND FUTURE WORK}
\pagebreak

\begin{center}
{\LARGE\textbf{CONCLUSIONS AND FUTURE WORK}}
\end{center}

% ─────────────────────────────────────────────────────────────────────────────

\section{Discussion of Results}
\label{sec:discussion_of_results}

It worked. Not perfectly, and not the way we originally imagined it would, but TrapPot ended up as a real, running system that we could demo, break, fix, and demo again.

Cowrie listened on ports 2222 and 2223. When we connected over SSH, it gave back a fake Linux shell and logged everything: the login attempt, the \texttt{whoami}, the \texttt{ls -la}, the \texttt{cat /etc/passwd}, and the \texttt{wget}. That is exactly what it is supposed to do, look like a vulnerable device and record whoever shows up \cite{9071014,9359711}. What surprised us is how fast the logs appeared. There was almost no lag between the command we typed and the JSON entry in Cowrie's output.

Zeek was the piece we were least confident about. It runs in Cowrie's network namespace (via \texttt{network\_mode: "service:cowrie"}) and has to observe traffic it does not own, so we disabled checksum validation with \texttt{-C} because containerized packet capture breaks without it. After the SSH test session, \texttt{conn.log} and \texttt{ssh.log} appeared with the correct fields. Twelve features. That is what the detector actually reads from those files, ports, protocol, service, duration, byte counts, packet counts, connection state, and the fact that Zeek produced them cleanly was the linchpin for everything downstream \cite{lanz2025optimizing,rodriguezgomez2025systematic}.

The AI detector read \texttt{conn.log}, extracted those features, loaded \texttt{trappot\_model.pkl} and \texttt{label\_encoder.pkl}, and printed live classification output. During testing it flagged the session as \texttt{Normal detected}, which is correct, since a single manual SSH session from localhost does not look like a brute-force attack. The result went to \texttt{detections.json}. Logstash picked that up, Elasticsearch indexed it, and Kibana showed it in the TrapPot Overview dashboard at \texttt{localhost:5601}. The full chain held.

The model itself hit 99.59\% accuracy on the IoT-23 hold-out set \cite{iot23dataset}. We cannot claim that number translates directly to live performance, the training data and the lab test environment are not the same thing, but it tells us the Random Forest classifier learned the five classes well from the features we gave it.

\subsection{Contributions}
\label{sub:contributions}

The main thing TrapPot contributes is a complete, portable workflow: one \texttt{docker compose up --build -d} and the whole stack starts, Cowrie, Zeek, the AI detector, Elasticsearch, Logstash, Kibana, and the setup containers \cite{9001068,lee2021classification}. That portability matters for an academic project. We were not installing seven tools on seven machines; we were running seven containers on one laptop.

The second contribution is the Zeek-as-bridge design. Cowrie sees attacker commands. Zeek sees network packets. The AI detector needs structured numbers, not raw text. By sitting Zeek inside Cowrie's network namespace and feeding its \texttt{conn.log} directly to the classifier, we got network-level evidence from honeypot-level traffic, without writing a custom parser \cite{franco2021survey,lanz2025optimizing}.

The third contribution is the training pipeline itself. The first three runs on the raw IoT-23 data came in around 72\%. The problem was class imbalance (brute-force records outnumbered malware download records by roughly 5:1). After we applied SMOTE oversampling and rebalanced the training set, accuracy jumped to 99.59\% on the 20\% hold-out \cite{saied2023comparative,karim2024efficient}. That lesson, that imbalanced data produces a model that ignores rare classes, is worth knowing.

\section{Limitations}
\label{sec:limitations}

\begin{itemize}

\item \textbf{We never exposed TrapPot to real internet traffic.} The whole system ran locally, behind Docker's bridge network. That was the right call ethically and safely, but it means we have no data on how the detector behaves against actual attackers, bots, script-kiddies, or automated Mirai-style scanners. The 99.59\% accuracy is real but it is a dataset number, not a deployment number \cite{iot23dataset,famera2025mirai}.

\item \textbf{The live test was thin.} One SSH session, five commands, one classification result. That is not nothing, but it is also not a stress test. We needed more varied traffic, repeated login failures, port scanning, wget with a real URL, to exercise the five classes against a live model.

\item \textbf{Twelve features from conn.log only.} The detector does not read Cowrie logs at all. It reads Zeek's network-level summary. That means command sequences, login-attempt counts, session duration from Cowrie's perspective, and file download indicators are all invisible to the classifier right now. Some of the most distinctive attack behavior happens at that layer \cite{9071014,rodriguezgomez2025systematic}.

\item \textbf{The ELK Stack is heavy.} Elasticsearch and Kibana together want at least 4~GB of RAM on the host. On a machine with less, containers restart during startup or throttle badly. The honeypot and detector are lightweight, the storage and visualization layer is not.

\end{itemize}

\section{Future Work}
\label{sec:future_work}

\subsection{Real Attack Traffic}
\label{subsec:realistic_attack_traffic}

The obvious next step is putting TrapPot in a real isolated lab network, not localhost, but a dedicated VLAN with proper firewall rules, and letting it collect traffic for 30 days \cite{lee2021classification,suareztangil2023stories}. We want to see what the detector does with actual port-22 scanners and Mirai-variant bots, not just \texttt{whoami} from a developer's terminal. That is the validation gap we could not close in this version.

\subsection{Combining Zeek and Cowrie Features}
\label{subsec:combined_features}

Right now the model sees twelve numbers from Zeek. It does not know that the attacker tried 847 password combinations, or ran \texttt{cat /etc/shadow}, or downloaded a MIPS binary. Merging Cowrie session features into the feature vector, login-attempt count, command rate, presence of known malicious commands, file transfer flag, would probably improve classification quality on the harder classes (command execution and malware download), and it would make the classifier aware of both sides of the attack \cite{9071014,dimauro2021supervised}.

\subsection{Streaming AI Results into Kibana}
\label{subsec:indexing_ai_results}

Currently the detector writes to \texttt{detections.json} and Logstash picks it up with some delay. A cleaner design would send predictions directly to Elasticsearch via its REST API, which would let Kibana show AI classification results in near-real time alongside the Cowrie and Zeek dashboards \cite{9001068,log2evt2025}. That is a straightforward change, maybe two hours of work, but we ran out of time.

\subsection{Model Comparison}
\label{subsec:additional_ml_models}

We chose Random Forest because Lee et al.\ got 96\% with it on nearly identical Cowrie data, and we had no reason to use something harder to explain \cite{lee2021classification,saied2023comparative}. But it would be worth comparing against XGBoost and a simple neural network on the same IoT-23 dataset. Whether the gap is one percent or ten percent matters for deciding whether complexity is worth it.

\subsection{Deployment Hardening}
\label{subsec:deployment_security}

The current setup exposes ports 2222, 2223, and 5601 with default credentials in a \texttt{.env} file. For anything beyond a classroom demo, that needs to change: proper secrets management, Kibana behind a reverse proxy with authentication, container resource limits, and a log retention policy. None of that is hard, it just was not the point of this version \cite{commey2025blockchain,franco2021survey}.

\section{Conclusion}
\label{sec:conclusion}

TrapPot set out to do one thing: build a working IoT honeypot that captures attacker-like traffic, classifies it with a machine learning model, and shows the result in a dashboard. The component checks in Chapter~\ref{sec:analysis_result} show that it does that.

The pipeline is Cowrie → Zeek → Random Forest → ELK. Every link was verified. The model hit 99.59\% on IoT-23. The dashboard loads at \texttt{localhost:5601}. And the whole thing starts with a single command, \texttt{docker compose up --build -d}, which means anyone on the team can reproduce it without re-reading 40 pages of setup instructions.

What it is not: a production system, a validated deployment, or a substitute for real-world traffic testing. Those are the gaps described in Sections~\ref{sec:limitations} and~\ref{sec:future_work}. But as a Docker-based proof of concept that classifies five attack classes and reached 99.59\% accuracy on the IoT-23 hold-out set \cite{iot23dataset}, it reached the target for this phase.
